\documentclass[12pt]{article}
\usepackage[T1]{fontenc}
\usepackage[utf8]{inputenc}
\usepackage{lmodern}
\usepackage{textcomp}
\usepackage{enumitem}
\usepackage{geometry}
\geometry{margin=1in}
\begin{document}
\begin{center}
Answer Key - Geography - Graduate Level Assessment 24 \
\small (Answers are ordered exactly as in the question paper.)
\end{center}

\section*{Multiple Choice}
\begin{enumerate}[label=\arabic*.]
\item \textbf{Answer:} A
\\ \textit{Options:} A) Lower respiratory infections; B) Preterm birth complications; C) Diarrheal diseases; D) Malaria
\\ \textit{Note:} Lower respiratory infections are the leading cause of mortality in children under five years old due to their high prevalence in low-income regions, where access to healthcare and preventive measures is limited, contributing to significant morbidity and mortality rates in this age group.
\item \textbf{Answer:} D
\\ \textit{Options:} A) Bipolar disorder; B) Schizophrenia; C) Alcohol use disorders; D) Anxiety disorders
\\ \textit{Note:} Anxiety disorders are the most prevalent mental disorders globally, affecting a significant portion of the population due to their high incidence and widespread impact on daily functioning and quality of life.
\item \textbf{Answer:} A
\\ \textit{Options:} A) by 4 fold; B) by 8 fold; C) by 16 fold; D) by 32 fold
\\ \textit{Note:} The correct answer is based on the significant economic growth India experienced during this period, which, when adjusted for inflation and purchasing power parity (PPP), indicates that the GDP per capita rose approximately fourfold due to factors such as liberalization policies, increased industrialization, and foreign investment.
\item \textbf{Answer:} C
\\ \textit{Options:} A) 4 billion; B) 3 billion; C) 2 billion; D) 1 billion
\\ \textit{Note:} The correct answer is based on the United Nations' projection that the global population of children aged 0 to 15 will remain approximately stable at around 2 billion by the year 2100 due to declining birth rates and demographic changes.
\item \textbf{Answer:} B
\\ \textit{Options:} A) 36\%; B) 56\%; C) 76\%; D) 96\%
\\ \textit{Note:} The correct answer of 56\% reflects historical data indicating that during the mid-20 th century, significant portions of the global population, particularly in developing regions, lacked access to education and literacy resources, resulting in a relatively low literacy rate.
\end{enumerate}

\section*{True / False}
\begin{enumerate}[label=\arabic*.]
\setcounter{enumi}{5}
\item \textbf{Answer:} False
\\ \textit{Options:} A) True; B) False
\\ \textit{Note:} The correct answer is: Benjamin Mkapa National Stadium
\item \textbf{Answer:} True
\\ \textit{Options:} A) True; B) False
\\ \textit{Note:} According to the passage: 7 July 1937
\item \textbf{Answer:} False
\\ \textit{Options:} A) True; B) False
\\ \textit{Note:} The correct answer is: Oporto region
\item \textbf{Answer:} True
\\ \textit{Options:} A) True; B) False
\\ \textit{Note:} According to the passage: Jan Ekier
\item \textbf{Answer:} False
\\ \textit{Options:} A) True; B) False
\\ \textit{Note:} The correct answer is: 1889 Universal Exposition
\end{enumerate}

\section*{Long Form Response}
\begin{enumerate}[label=\arabic*.]
\setcounter{enumi}{10}
\item \textbf{Answer:} two years
\\ \textit{Note:} Expected answer: two years
\item \textbf{Answer:} 6800
\\ \textit{Note:} Expected answer: 6800
\end{enumerate}

\vfill
\noindent\textit{End of Answer Key}
\end{document}